\documentclass[a4paper,12pt]{article}
\usepackage[utf8]{inputenc}
\usepackage[T1]{fontenc}
\usepackage[ngerman]{babel}
\usepackage{amsmath}
\usepackage{amssymb}
\usepackage{enumitem}
\usepackage{geometry}
\geometry{a4paper, margin=2.5cm}

\title{Einsendeaufgaben -- Lektion 4\\[0.5em]
\large Modul 61111: Mathematische Grundlagen}
\author{}
\date{}

\begin{document}
\maketitle

\section*{Aufgabe 4.2}

\begin{enumerate}[label=\alph*)]

    \item \underline{$\sup M_1 = 1$}

    Beweis:

    \begin{enumerate}[label=\arabic*.]
        \item 1 ist eine obere Schranke, da
        \[
            \frac{n-1}{n+1} < 1 \Leftrightarrow n-1 < n+1 \Leftrightarrow -1 < 1.
        \]

        \item Es gilt.
        \[
            \lim_{n \to \infty} \frac{n-1}{n+1}
            = \lim_{n \to \infty} \frac{1 - \frac{1}{n}}{1 + \frac{1}{n}} = 1.
        \]
        Dann sei $\varepsilon > 0$ beliebig. Wegen $\dfrac{1-\frac{1}{n}}{1+\frac{1}{n}} \to 1$,
        existiert ein $n \in \mathbb{N}$, sodass
        \[
            \left| \frac{1-\frac{1}{n}}{1+\frac{1}{n}} - 1 \right| < \varepsilon
        \]
        Da 1 obere Schranke ist, folgt
        \[
            \left| \frac{1-\frac{1}{n}}{1+\frac{1}{n}} - 1 \right|
            = |-1| \left| 1 - \frac{1-\frac{1}{n}}{1+\frac{1}{n}} \right|
            = 1 - \frac{1-\frac{1}{n}}{1+\frac{1}{n}} < \varepsilon.
        \]
        Also gilt
        \[
            1 - \varepsilon < \frac{1-\frac{1}{n}}{1+\frac{1}{n}}
        \]
    \end{enumerate}

    \underline{max $M_1$:}
    existiert nicht, da $\sup M_1 \notin M_1$

    \underline{$\min M_1 = 0$}

    Beweis:

    Da $n \geq 1$ für alle $n \in \mathbb{N}$ gilt, folgt
    $0 \leq n-1$ und damit $0 \leq \dfrac{n-1}{n+1}$.
    Und $0 \in M_1$, da für $n=1$
    \[
        0 = \frac{1-1}{1+1} = \frac{n-1}{n+1} \text{ gilt.}
    \]

    \underline{$\inf M_1 = 0$}

    Beweis:

    folgt aus $\min M_1 = \inf M_1$, da Minimum existiert.

    \item $M_2$ besitzt keine untere Schranke.

    Beweis:

    Wir definieren eine Folge $(a_n)$ mit
    \[
        a_n = \frac{\left(-1+\frac{1}{n}\right)}{1 + \left(-1+\frac{1}{n}\right)}
        = \frac{-1+\frac{1}{n}}{\frac{1}{n}} = -n+1
    \]
    Die Folge $(a_n)$ besitzt keine untere Schranke.
    Da jedes $a_n \in M_2$, existiert keine untere Schranke für $M_2$.

    Also existiert kein $\min M_2$ und kein $\inf M_2$.

    Allerdings existiert eine obere Schranke; nämlich die 1:

    Als $x < x+1$ für alle $x \in (-1, \infty)$
    folgt $\dfrac{x}{x+1} < 1$.

    Die 1 ist sogar das Supremum, da
    die Folge $(a_n)$ mit $a_n = \dfrac{n}{n+1}$ gegen 1
    konvergiert und $a_n \in M_2$.
    $\sup M_2$ liegt nicht in $M_2$, da aus der Annahme
    \[
        \frac{x}{1+x} = 1
    \]
    der Widerspruch $x = 1+x$
    folgt.
    Also existiert kein Maximum.

\end{enumerate}

\end{document}
